\chapter{Hogyan gondolkodjunk egy RWLang alkalmazásról?}

Mielőtt telepítünk, route-ot írunk vagy adatbázist kötünk az alkalmazáshoz, érdemes tisztázni egy szemléleti kérdést. Sok webes technológiában természetes kiindulópont az oldal, az endpoint vagy a controller: \emph{``kell egy \code{/products} URL, írjunk hozzá handlert''}. RWLangban hasznosabb előbb azt megkérdezni: \emph{``milyen üzleti adat létezik, milyen alakban, milyen műveletek végezhetők rajta, és milyen határok között?''}

Ez nem akadémiai különbség. A nyelv több biztonsági és helyességi garanciája abból származik, hogy az alkalmazás fontos adatalakjai, bemenetei, lekérdezései és erőforrás-hozzáférései explicit szerződések.

\section{Először a fogalmak, utána az oldalak}

Vegyünk egy egyszerű terméklistát. A feladatot lehetne így kezdeni:

\begin{quote}
Kell egy GET oldal, egy POST endpoint és egy SQL INSERT.
\end{quote}

RWLangban célszerűbb így gondolkodni:

\begin{enumerate}
  \item Van egy üzleti fogalom: \emph{Product}.
  \item A programban ennek ismert adatalakja van: \code{id}, \code{name}, \code{price}.
  \item Vannak olvasó és módosító műveletek: listázás és létrehozás.
  \item A módosítás adatbázis-tranzakciót igényel.
  \item A HTTP réteg typed inputot ad a műveletnek, és policy-t rendel hozzá.
  \item A megjelenítés már az így kapott typed adatot rendereli HTML-be vagy JSON-ba.
\end{enumerate}

A gondolkodás iránya tehát nagyjából:

\begin{lstlisting}
uzleti fogalom
    -> typed adatalak
        -> muveletek
            -> capability-k es policy-k
                -> HTTP route
                    -> HTML / JSON
\end{lstlisting}

Nem minden alkalmazásnál kell ezt papíron végigírni. A sorrend mégis segít elkerülni, hogy az üzleti szabályok véletlenül route-okba, template-ekbe vagy kézzel összefűzött SQL-be szóródjanak szét.

\section{Mit jelent a \code{model}?}

A könyvben két közeli fogalommal találkozunk:

\begin{itemize}
  \item \textbf{adatmodell vagy domain modell}: annak emberi leírása, milyen üzleti fogalmak és kapcsolatok léteznek az alkalmazásban;
  \item \textbf{\code{model} deklaráció}: az RWLang konkrét, fordító által ismert adatalakja.
\end{itemize}

Például:

\begin{lstlisting}[language=RWLang]
model Product {
    id: Int
    name: String
    price: Int
}
\end{lstlisting}

Ez azt mondja a fordítónak, hogy egy \code{Product} rekordnak pontosan ilyen mezői és típusai vannak. A typed SQL query eredményalakja ehhez a szerződéshez igazodik, és ugyanez az érték később HTML- vagy JSON-kimenetben is használható.

\subsection*{A \code{model} nem automatikus adatbázisséma}

A fenti deklaráció önmagában nem hoz létre \code{products} táblát. Az adatbázissémát migration hozza létre, például:

\begin{lstlisting}[language=SQL]
CREATE TABLE products (
    id INTEGER PRIMARY KEY,
    name TEXT NOT NULL,
    price INTEGER NOT NULL
);
\end{lstlisting}

A két réteg kapcsolatban áll, de nem ugyanaz. Ez tudatos szétválasztás:

\begin{itemize}
  \item a migration az adatbázis tartós sémájának életciklusa;
  \item a \code{model} az alkalmazáskód typed adatszerződése;
  \item a query mondja meg explicit módon, hogyan lesz adatbázissorokból model rekord.
\end{itemize}

Ezért RWLangban nincs szükség arra a feltételezésre, hogy minden model automatikusan egy ORM-entitás vagy egy azonos nevű tábla.

\subsection*{A \code{model} nem mutable objektum}

A \code{Product} nem klasszikus objektumorientált osztály. Nincs mögötte rejtett \code{save()} metódus, implicit adatbázis-kapcsolat vagy mutable objektuméletciklus. Ha adatot akarunk módosítani, explicit műveletet és explicit erőforrást használunk.

Például az olvasás:

\begin{lstlisting}[language=RWLang]
query fn listProducts(db: Db) -> Result<List<Product>, DbError> sql {
    SELECT id, name, price
    FROM products
    ORDER BY id DESC
    LIMIT 50
}
\end{lstlisting}

A módosítás pedig külön, tranzakcióhoz kötve jelenik meg:

\begin{lstlisting}[language=RWLang]
query fn createProduct(
    tx: Transaction,
    name: String,
    price: Int
) -> Result<Void, DbError> sql {
    INSERT INTO products(name, price)
    VALUES (:name, :price)
}
\end{lstlisting}

A függvény aláírásából látszik, hogy melyik művelet csak olvas, és melyik kap módosításra alkalmas tranzakciót.

\section{A model a rétegek közti szerződés}

Érdemes a \code{model}-re nem pusztán ``adatbázisos struktúraként'', hanem a rétegek között mozgó typed szerződésként gondolni.

\begin{lstlisting}
Database row
    |  typed query
    v
Product model
    |-------------------|
    v                   v
HTML page            JSON API
    |
    v
felhasznalo
\end{lstlisting}

A fordító emiatt tud olyan hibákat korán megfogni, amelyeket dinamikus rendszerekben csak futás közben vennénk észre. Ha egy query \code{Product}-ot ígér, a kiválasztott/returned mezőknek a model shape-jéhez kell igazodniuk. Nem az a cél, hogy a program ``majd valahogy'' értelmezze az eredményt.

Ez a szemlélet később az object authorizationnél és a JSON API-nál is visszatér: a rendszer ismert, typed rekordokkal dolgozik ahelyett, hogy tetszőleges, futás közben változó objektumokat adna tovább.

\section{Az üzleti fogalomhoz műveletek tartoznak}

Egy model önmagában ritkán érdekes. A valódi alkalmazásban műveletek tartoznak hozzá:

\begin{itemize}
  \item \code{Product.list};
  \item \code{Product.byId};
  \item \code{Product.create};
  \item egy CMS-ben például \code{Article.publish};
  \item egy wikiben \code{WikiPage.update}.
\end{itemize}

Kis programban ezek lehetnek top-level függvények. Nagyobb programban a később bemutatott \code{object} domain namespace fogja őket egy üzleti fogalom köré rendezni.

Fontos különbség:

\begin{center}
\begin{tabularx}{\textwidth}{@{}lY@{}}
\toprule
Elem & Kérdés, amelyre válaszol \\
\midrule
\code{model} & Milyen alakú adat a \code{Product}? \\
\code{query fn} & Hogyan olvassuk vagy módosítjuk a tartós adatot? \\
\code{object} & Mely üzleti fogalomhoz tartoznak ezek a műveletek? \\
\code{page}/\code{action} & Hogyan kapcsolódik az üzleti művelet a webes handlerhez? \\
\code{route} & Milyen HTTP útvonalon, inputtal és policy-val érhető el? \\
\bottomrule
\end{tabularx}
\end{center}

\section{Capability: a hozzáférés is része a modellnek}

RWLangban nem csak az adat típusa explicit, hanem az is, milyen erőforráshoz fér hozzá egy művelet. Egy függvény nem ``valahonnan'' kap adatbázist, fájlrendszert vagy current usert.

\begin{lstlisting}[language=RWLang]
page fn home(ctx: PageContext, db: Db)
    -> Result<Html, PageError> {
    let products = listProducts(db)?;
    return Ok(html {
        <ul>
            @for product in products {
                <li>{{ product.name }}</li>
            }
        </ul>
    });
}
\end{lstlisting}

A \code{db: Db} az aláírásban látható. Fájlkezelésnél AppFs capability, módosító adatbázis-műveletnél Transaction, autentikált műveletnél pedig az auth policy és a hozzá tartozó typed kontextus válik relevánssá.

Ebből egy hasznos tervezési szabály következik:

\begin{quote}
Ha egy műveletnek érzékeny erőforrás kell, annak a kódból és a policy-ból látszania kell.
\end{quote}

\section{A route nem az üzleti modell}

A route fontos, de a rendszer külső határa. Például:

\begin{lstlisting}[language=RWLang]
route create POST "/products"
    form name<String> price<Int>
    validate name length 1 100 price range 0 100000000
    => create;
\end{lstlisting}

Ez azt mondja meg, hogyan léphet be egy HTTP kérés a rendszerbe: melyik metóduson és pathon, milyen typed inputtal, milyen validációval és később milyen auth/cache/rate-limit policy-val.

Az üzleti fogalom azonban nem a \code{/products} URL. Ugyanaz a \code{Product} később HTML oldalon, JSON API-ban, admin felületen vagy más route mögött is használható.

Ezért nagyobb alkalmazásban különösen hasznos a következő határ:

\begin{lstlisting}
domain/model     : mi az adat es milyen muvelet ertelmes rajta?
HTTP route       : ki, hogyan es milyen policy-val hivhatja?
runtime config   : milyen infrastruktura-es eroforras-hatarok kozott futhat?
\end{lstlisting}

\section{Egy egyszerű tervezési módszer}

Amikor új feature-t kezdesz, a következő öt kérdés jó kiindulópont:

\begin{enumerate}
  \item \textbf{Mi az üzleti főnév?} -- Product, Article, Invoice, WikiPage.
  \item \textbf{Milyen typed adatra van ténylegesen szükség?} -- ne automatikusan a teljes adatbázissort másold át.
  \item \textbf{Mely műveletek olvasnak, és melyek módosítanak?}
  \item \textbf{Milyen capability kell hozzájuk?} -- Db, Transaction, AppFs, auth, később hálózati egress.
  \item \textbf{Milyen HTTP/policy felületen tesszük elérhetővé?} -- route, input, validáció, auth, cache, rate limit.
\end{enumerate}

Egy egyszerű CRUD esetén ez pár perc gondolkodás. Egy CMS-nél vagy ügyviteli rendszernél viszont ez akadályozza meg, hogy az alkalmazás néhány hónap után route-ok, SQL-ek és implicit jogosultsági feltételek nehezen követhető halmazává váljon.

\section{Mit vigyél tovább a következő fejezetekbe?}

Nem kell minden alkalmazást ``domain-driven design'' terminológiával megtervezni, és nem kell minden hárommezős rekordból külön object hierarchy-t építeni. A könyvben használt modell-szemlélet ennél egyszerűbb:

\begin{itemize}
  \item nevezd meg a fontos üzleti adatot;
  \item add meg a typed alakját;
  \item az adatmozgást explicit query-kkel írd le;
  \item a módosítást explicit tranzakcióhoz kösd;
  \item az erőforrás-hozzáférést capability-ként tedd láthatóvá;
  \item a HTTP-t és a security policy-t külső határként kezeld.
\end{itemize}

A következő fejezetekben először lefordítjuk, telepítjük és csak a szükséges minimumon konfiguráljuk a rendszert, majd az A--Z példában ugyanezt a szemléletet egy működő \code{Product} alkalmazáson végigkövetjük. Utána a könyv minden egyes részét külön mélyíti el.
